\section*{Paper Structure Summary}

\textbf{Working Title}

\begin{center}
\textit{Tokenizing Labor: A Probabilistic Framework for Estimating Compute, Energy, and Infrastructure Requirements of AI Software Engineering Substitution}
\end{center}

\vspace{0.5em}

\section{Introduction}

Artificial intelligence (AI) systems are increasingly capable of performing software engineering tasks, leading to growing discussion of AI-driven labor substitution. At the same time, AI has emerged as a significant driver of compute infrastructure deployment, data-center expansion, and electricity demand. Existing studies primarily estimate AI energy demand from projected compute growth, GPU deployment, or data-center capacity expansion. However, little work has examined how labor substitution itself may generate downstream demand for tokens, compute resources, electricity, and infrastructure.

Software engineering provides a particularly suitable case study because it represents a large global workforce, operates almost entirely within digital environments, possesses rich benchmark ecosystems, and exhibits rapidly improving AI capability. Furthermore, software engineering work can be naturally translated into token-based AI workloads.

This paper addresses the following research questions:

\begin{enumerate}
\item How can software-engineering labor substitution be translated into compute, energy, and infrastructure demand?
\item What infrastructure requirements emerge under different AI adoption scenarios?
\item Which uncertainties dominate these estimates?
\item What implications do these findings have for energy systems and infrastructure planning?
\end{enumerate}

The primary contributions of this work are:

\begin{enumerate}
\item A labor-to-infrastructure framework linking labor substitution to compute and energy demand.
\item The introduction of \textbf{Tokenized Labor Intensity (TLI)} as a measure of token demand per unit of labor.
\item The introduction of \textbf{Levelized Cost of Engineering Substitution (LCES)} as a metric for comparing AI infrastructure costs against software labor costs.
\item A probabilistic infrastructure assessment framework based on uncertainty propagation and Monte Carlo simulation.
\end{enumerate}

Figure~1 presents the overall framework linking labor, tokens, compute, energy, infrastructure, and economics.

\section{Conceptual Framework}

The central conceptual contribution of this paper is the proposition that tokens serve as the bridge between labor economics and energy economics.

We define \textbf{Tokenized Labor Intensity (TLI)} as

[
TLI =
\frac{\text{Tokens}}
{\text{Human Labor}}
]

which may be expressed as:

[
TLI =
\frac{\text{Tokens}}
{\text{FTE-day}}
]

or

[
TLI =
\frac{\text{Tokens}}
{\text{Labor Hour}}
]

The framework follows the chain

[
\text{Labor}
\rightarrow
\text{Tokens}
\rightarrow
\text{Compute}
\rightarrow
\text{Energy}
\rightarrow
\text{Infrastructure}
\rightarrow
\text{Economics}
]

Although software engineering is used as the case study, the framework may be generalized to other knowledge-work occupations such as legal services, accounting, customer support, consulting, and research.

\section{Methodology}

\subsection{Deterministic Accounting Model}

The deterministic model consists of five linked modules.

\paragraph{Labor Module}

The number of substituted software-engineering workers is

[
N_{\mathrm{replaced}}
=====================

N_{\mathrm{eng}}
\cdot
s
]

where

\begin{itemize}
\item $N_{\mathrm{eng}}$ is the total software-engineering workforce,
\item $s$ is the FTE-equivalent substitution rate.
\end{itemize}

\paragraph{Token Module}

Total daily token demand is

[
D
=

N_{\mathrm{replaced}}
\cdot
T_{\mathrm{base}}
\cdot
A_{\mathrm{agent}}
]

where

\begin{itemize}
\item $T_{\mathrm{base}}$ is baseline token demand per substituted FTE,
\item $A_{\mathrm{agent}}$ is an amplification factor representing context loading, retries, validation, and tool usage.
\end{itemize}

\paragraph{Compute Module}

Required throughput is

[
R
=

\frac{D}{86400}
]

and required GPU count is

[
G
=

\frac{R}
{\theta_{\mathrm{GPU}} U}
]

where

\begin{itemize}
\item $\theta_{\mathrm{GPU}}$ is effective GPU throughput,
\item $U$ is utilization.
\end{itemize}

\paragraph{Energy Module}

IT load is

[
P_{\mathrm{IT}}
===============

G
\cdot
P_{\mathrm{GPU}}
\cdot
\alpha_{\mathrm{system}}
]

Facility load is

[
P_{\mathrm{facility}}
=====================

P_{\mathrm{IT}}
\cdot
PUE
]

Annual electricity demand is

[
E_{\mathrm{annual}}
===================

P_{\mathrm{facility}}
\cdot
8760
]

where

\begin{itemize}
\item $P_{\mathrm{GPU}}$ is average GPU power,
\item $\alpha_{\mathrm{system}}$ is server and networking overhead,
\item $PUE$ is data-center power usage effectiveness.
\end{itemize}

\paragraph{Economic Module}

Annual energy cost is

[
C_{\mathrm{energy}}
===================

E_{\mathrm{annual}}
\cdot
C_{\mathrm{MWh}}
]

Infrastructure capital cost is

[
C_{\mathrm{capex}}
==================

G
\cdot
C_{\mathrm{GPU}}
+
P_{\mathrm{facility}}
\cdot
C_{\mathrm{DC}}
]

Annualized infrastructure cost is

[
C_{\mathrm{annual}}
===================

CRF
\cdot
C_{\mathrm{capex}}
]

where

[
CRF
===

\frac{i(1+i)^n}
{(1+i)^n-1}
]

The \textbf{Levelized Cost of Engineering Substitution (LCES)} is defined as

[
LCES
====

\frac{
C_{\mathrm{annual}}
+
C_{\mathrm{energy}}
+
C_{\mathrm{O&M}}
}
{
N_{\mathrm{replaced}}
}
]

The annual labor-equivalent value is

[
V_{\mathrm{labor}}
==================

N_{\mathrm{replaced}}
\cdot
C_{\mathrm{FTE}}
]

and the AI-to-labor cost ratio is

[
\rho
====

\frac{
LCES
}{
C_{\mathrm{FTE}}
}
]

\subsection{Scenario Design}

Three scenarios are evaluated:

\begin{enumerate}
\item AI-Assisted Engineering
\item Agentic Engineering
\item Autonomous Engineering Organization
\end{enumerate}

Each scenario is characterized by distinct assumptions regarding substitution rates, token intensity, retry behavior, throughput, and utilization.

\subsection{Uncertainty Characterization}

Model parameters are represented using probability distributions.

Examples include:

\begin{itemize}
\item Workforce size
\item Tokenized Labor Intensity
\item Agentic amplification factor
\item GPU throughput
\item Utilization
\item GPU power
\item PUE
\item Electricity price
\item Infrastructure costs
\item Labor compensation
\end{itemize}

Table~1 summarizes all parameter distributions and uncertainty ranges.

\subsection{Probabilistic Assessment}

A Latin Hypercube Monte Carlo framework is employed to propagate uncertainty through the deterministic model.

For each simulation, the framework generates estimates of:

\begin{itemize}
\item Daily token demand
\item Required GPU count
\item Facility power demand
\item Annual electricity consumption
\item Infrastructure cost
\item LCES
\end{itemize}

\section{Deterministic Baseline Assessment}

A baseline Agentic Engineering scenario is evaluated using central parameter values.

Outputs include:

\begin{itemize}
\item Token demand
\item GPU fleet size
\item Facility load
\item Annual electricity demand
\item Capital investment
\item LCES
\end{itemize}

The purpose of this section is to establish order-of-magnitude intuition before introducing uncertainty analysis.

\section{Probabilistic Results}

Monte Carlo simulations are used to generate output distributions for:

\begin{itemize}
\item Token demand
\item GPU count
\item Facility power
\item Annual electricity demand
\item LCES
\end{itemize}

Results are reported as:

\begin{itemize}
\item Mean
\item Median
\item 5th percentile
\item 95th percentile
\end{itemize}

Scenario comparisons are performed across all three AI adoption pathways.

Threshold analyses evaluate probabilities such as

[
P(P_{\mathrm{facility}} > 100\ \mathrm{MW})
]

[
P(P_{\mathrm{facility}} > 500\ \mathrm{MW})
]

and

[
P(LCES < 0.1 \times C_{\mathrm{FTE}})
]

\section{Sensitivity Analysis}

Sensitivity analysis identifies the dominant drivers of infrastructure and energy demand.

A tornado analysis is first conducted using one-at-a-time parameter variation.

Global sensitivity analysis is then performed using Spearman rank correlations and, optionally, Sobol indices.

Particular attention is given to the influence of:

\begin{itemize}
\item Tokenized Labor Intensity
\item Agentic amplification
\item GPU throughput
\item Utilization
\item PUE
\item Electricity price
\end{itemize}

The objective is to determine whether AI-system variables or conventional energy-system variables dominate uncertainty.

\section{Discussion}

The discussion focuses on five themes.

First, the paper introduces tokenized labor as a new demand-side framework for estimating AI infrastructure growth.

Second, it examines implications for utilities, data-center developers, cloud providers, and grid planners.

Third, it analyzes why AI-system variables may exert greater influence on infrastructure requirements than traditional energy-system variables such as PUE or electricity price.

Fourth, it discusses rebound effects, whereby productivity gains may increase total software production and therefore increase AI demand.

Finally, it considers the broader applicability of the framework to other knowledge-work occupations.

\section{Limitations}

Key limitations include:

\begin{itemize}
\item Benchmark performance does not necessarily imply labor substitution.
\item Token-demand estimates remain uncertain.
\item Hardware efficiency may evolve rapidly.
\item Organizational adoption remains difficult to forecast.
\item Rebound effects are not fully modeled.
\end{itemize}

\section{Conclusions}

This paper proposes a labor-to-infrastructure framework that links software-engineering labor substitution to token demand, compute requirements, energy consumption, and economic outcomes.

By introducing Tokenized Labor Intensity (TLI) and Levelized Cost of Engineering Substitution (LCES), the framework establishes a quantitative bridge between labor economics and energy economics.

The results demonstrate how AI-driven labor substitution can be translated into infrastructure demand and provide a probabilistic basis for evaluating future AI-related energy-system impacts.

